% Diagrama 2. Flujo de la secuencia de remoción, en TikZ.
% Equivale a D2_flujo_remocion.mmd. (P) = propuesto, no fijado en la tesis.
% Actualizado el 19 oct 2026 a la tesis del 15 oct: modulación del ventilador, temperatura de la impresora,
% L12-P con posición, tiempos máximos de la Tabla 15 y tres intentos.
% 21 oct 2026 (marcas de Diego): la autorización del usuario solo se pide si la confianza del agente es baja;
% "botón de emergencia" en lugar de "seta"; referencias por nombre; rótulos separados de las flechas.
% Compilar: pdflatex D2_flujo_remocion_tikz.tex
\documentclass[border=4pt]{standalone}
\usepackage[T1]{fontenc}
\usepackage[utf8]{inputenc}
\usepackage[spanish,es-noshorthands]{babel}
\usepackage{times}
\usepackage{amssymb}
\usepackage{tikz}
% Estilos comunes de los diagramas de validación (TikZ).
% Blanco y negro, Times a 8 pt, sin colores, como la Figura 22 de la tesis.
\usetikzlibrary{positioning,calc,shapes.geometric}
\tikzset{
  base/.style   = {font=\fontsize{8}{9.6}\selectfont},
  proc/.style   = {draw, fill=white, align=center, text width=3.0cm,
                   minimum height=0.9cm, inner sep=3pt,
                   execute at begin node={\hyphenpenalty=10000\relax}},
  propp/.style  = {proc, dashed},
  negra/.style  = {proc, fill=black!12},
  falla/.style  = {proc, fill=black!7},
  dec/.style    = {draw, fill=white, diamond, aspect=1.9, align=center,
                   inner sep=0.5pt, text width=1.95cm,
                   execute at begin node={\hyphenpenalty=10000\relax}},
  term/.style   = {draw, fill=black!4, rounded corners=0.42cm, align=center,
                   text width=2.7cm, minimum height=0.84cm, inner sep=3pt,
                   execute at begin node={\hyphenpenalty=10000\relax}},
  con/.style    = {draw, circle, fill=white, minimum size=0.6cm, inner sep=0pt},
  nota/.style   = {draw, dashed, fill=white, align=center, text width=4.6cm,
                   inner sep=3pt, font=\fontsize{7.2}{8.6}\selectfont},
  marco/.style  = {draw, rounded corners=5pt, black!55, thin},
  tmarco/.style = {font=\fontsize{8}{9.6}\selectfont\itshape, anchor=north west},
  fl/.style     = {->, >=latex, semithick},
  ln/.style     = {semithick},
  et/.style     = {font=\fontsize{6.8}{8}\selectfont, inner sep=1.2pt, fill=white},
  union/.style  = {circle, fill=black, inner sep=0pt, minimum size=3pt}
}


\begin{document}
\begin{tikzpicture}[base]
% ======================= parte 1: preparación =======================
% 22 oct 2026: la consulta al agente y, si la confianza es baja, la autorización del usuario se hacen al terminar
% la impresión y antes del enfriamiento, como en el pseudocódigo, la Tabla de interfaz y la implementación parcial.
\node[term] (A)  at (3.0,0)     {Fin de impresión\\(estado por red local)};
\node[dec, text width=1.7cm]  (E)  at (3.0,-1.95) {¿Pieza presente? (captura)};
\node[term, text width=1.6cm] (L0) at (-0.6,-1.95) {Placa libre: fin sin flexión};
\node[dec, text width=1.6cm]  (CF) at (3.0,-4.3) {¿Confianza del agente alta?};
\node[dec, text width=1.6cm]  (G)  at (3.0,-6.6) {¿Autoriza el usuario? (R-SEG-01)};
\node[falla, text width=1.9cm] (H1) at (6.85,-6.6) {Espera de intervención humana};
\node[dec]  (B)  at (3.0,-9.1) {¿T > 38.7 °C?};
\node[proc] (C)  at (3.0,-10.8) {\textbf{Etapa 1}\\ventilador principal al 100\,\% por modulación (D27), nominal 3.3 min};
\node[dec]  (C1) at (3.0,-12.7) {¿T $\leq$ 38.7 °C? (impresora)};
\node[falla, text width=1.9cm] (FT) at (6.85,-14.35) {Falla del enfriamiento: aviso al usuario};
\node[con]  (X1) at (3.0,-14.95) {1};

\draw[fl] (A) -- (E);
\draw[fl] (E.west) -- node[et, above=2pt, pos=0.5] {no} (L0.east);
\draw[fl] (E) -- node[et, right=1pt] {sí} (CF);
\draw[fl] (CF) -- node[et, right=1pt] {no} (G);
\draw[ln] (CF.west) -- node[et, above=2pt, pos=0.5] {sí} (0.0,-4.3) -- (0.0,-7.9) -- (3.0,-7.9);
\node[union] at (3.0,-7.9) {};
\draw[fl] (G) -- node[et, right=1pt, pos=0.35] {sí} (B);
\draw[fl] (G.east) -- node[et, above=2pt, align=center, pos=0.5] {no, o 30 min\\sin respuesta} (H1.west);
\draw[fl] (B) -- node[et, right=1pt] {sí} (C);
\draw[fl] (C) -- (C1);
\draw[fl] (C1.west) -- (1.25,-12.7) -- node[et, left=1pt, align=right, pos=0.5] {no, antes\\de 6 min} (1.25,-10.8) -- (C.west);
\draw[fl] (C1.east) -- node[et, above=2pt, align=center, pos=0.45] {no, pasados 6 min,\\o sin giro} (6.85,-12.7) -- (FT.north);
\draw[fl] (C1) -- node[et, right=1pt, pos=0.3, align=left] {sí, o sin dato\\y 6 min cumplidos} (X1);
\draw[ln] (B.west) -- node[et, above=2pt, pos=0.6] {no} (-0.6,-9.1) -- (-0.6,-14.25) -- (3.0,-14.25);
\node[union] at (3.0,-14.25) {};

\draw[marco] (-1.75,1.0) rectangle (8.1,-15.5);
\node[tmarco] at (-1.65,0.95) {Parte 1: preparación};

% ======================= parte 2: ejecución =======================
\node[con]  (X2) at (12.9,0)     {1};
\node[proc, text width=3.3cm] (I)  at (12.9,-1.75) {\textbf{Etapa 2}\\extensión hasta 39.4 mm leídos en el potenciómetro (D34), 7.8 s};
\node[dec, text width=1.8cm]  (I1) at (12.9,-3.85) {¿39.4 mm antes de 12 s y sin bloqueo?};
\node[falla, text width=1.45cm] (BQ) at (16.3,-3.85) {Bloqueo: retracción y aviso};
\node[proc, text width=3.3cm] (J)  at (12.9,-5.95) {\textbf{Etapa 3}\\retracción del actuador y barrido del cabezal a tres cotas, 30 s};
\node[dec, text width=1.6cm]  (J1) at (12.9,-7.95) {¿Cabezal en reposo antes de 60 s?};
\node[proc, text width=3.3cm] (K)  at (12.9,-9.95) {\textbf{Etapa 4}\\iluminación, captura y consulta al agente remoto, 5 s};
\node[dec]  (K1) at (12.9,-11.65) {¿Placa libre?};
\node[term] (OK) at (12.9,-13.45) {Aviso al usuario: equipo disponible};
\node[dec, text width=1.35cm] (R) at (9.85,-11.65) {¿Intento $<$ 3?};
\node[term, text width=1.9cm, dashed] (SRT) at (15.75,-9.95) {Estado Sin red del diagrama de estados};
\node[falla, text width=1.9cm] (H2) at (15.75,-13.45) {Espera de intervención humana};

\draw[fl] (X2) -- (I);
\draw[fl] (I) -- (I1);
\draw[fl] (I1) -- node[et, right=1pt] {sí} (J);
\draw[fl] (I1.east) -- node[et, above=2pt] {no} (BQ.west);
\draw[fl] (BQ.east) -- (17.3,-3.85) -- (17.3,-13.45) -- (H2.east);
\draw[fl] (J) -- (J1);
\draw[fl] (J1) -- node[et, right=1pt] {sí} (K);
\draw[fl] (K) -- (K1);
\draw[fl] (K1) -- node[et, right=1pt] {sí} (OK);
\draw[fl] (K1.west) -- node[et, above=2pt] {no} (R.east);
\draw[fl] (R.north) -- node[et, left=1pt, pos=0.5] {sí} (9.85,-1.75) -- (I.west);
\draw[fl] (J1.east) -- node[et, above=2pt, align=center, pos=0.45] {no, o\\sin red} (15.75,-7.95) -- (SRT.north);
\draw[fl] (K1.east) -- node[et, above=2pt, align=center, pos=0.45] {sin respuesta\\en 30 s} (15.75,-11.65) -- (SRT.south);
\draw[fl] (R.south) -- node[et, left=1pt, pos=0.35] {no} (9.85,-14.35) -- (15.75,-14.35) -- (H2.south);

\draw[marco] (8.3,1.0) rectangle (17.6,-15.5);
\node[tmarco] at (8.4,0.95) {Parte 2: ejecución};

\node[nota, text width=12.5cm] at (7.93,-16.1) {En cualquier punto: botón de emergencia pulsado, corte de +24V\_ACT (parada de emergencia del diagrama de estados).};
\end{tikzpicture}
\end{document}
